\section*{Capítulo \ref{ch:b1:clastic-sediments} --- Sedimentos clásticos y sus ambientes}

\begin{solution}{exo:b1:clastic-sediments:1}
$\varphi = -\log_2 d$: \qty{0.18}{mm} da $2.47$, una arena fina (2 a 3);
\qty{0.75}{mm} da $0.42$, una arena gruesa (0 a 1); \qty{3}{mm} da
$-1.58$, un gránulo ($-1$ a $-2$); \qty{0.012}{mm} da $6.38$, un limo fino (6
a 7); \qty{150}{mm} da $-7.23$, un canto rodado ($-6$ a $-8$).
\end{solution}

\begin{solution}{exo:b1:clastic-sediments:2}
$M_z = (1.05 + 1.71 + 2.43)/3 = 1.73\,\varphi$, es decir $2^{-1.73} =
\qty{0.30}{mm}$, una arena media. $\sigma_I = (2.43 - 1.05)/4 + (2.98 - 0.62)/6.6
= 0.345 + 0.358 = 0.70\,\varphi$: moderadamente bien clasificada (0.50 a 0.71), justo
por debajo del límite de moderadamente clasificada.
\end{solution}

\begin{solution}{exo:b1:clastic-sediments:3}
$q_b = (1 - 0.4) \times 0.02 \times (\num{0.01}/60)/2 = \qty{1.0e-6}{m^2.s^{-1}}$
de granos sólidos. En un día, $\num{1.0e-6} \times 86400 = \qty{0.0864}{m^3}$ por
metro de cresta, o $0.0864 \times 2650 = \qty{229}{kg}$ de cuarzo.
\end{solution}

\begin{solution}{exo:b1:clastic-sediments:4}
(a)~50 nudos son $50 \times 1852/3600 = \qty{25.7}{m.s^{-1}}$; 12 nudos,
\qty{6.17}{m.s^{-1}}. (b)~$\qty{100000}{m}/\qty{6.17}{m.s^{-1}} =
\qty{16200}{s}$, unas \qty{4.5}{h}. (c)~480 millas son \qty{772}{km}, 400 millas
\qty{644}{km}.
\end{solution}

\begin{solution}{exo:b1:clastic-sediments:5}
(a)~$(\varphi_{84} - \varphi_{16})/2 = 2 \times 0.9945\,\sigma/2 =
0.9945\,\sigma$, y $(\varphi_{95} - \varphi_5)/3.3 = 2 \times
1.6449\,\sigma/3.3 = 0.9969\,\sigma$: ambas estiman $\sigma$ con menos de un tercio de un
por ciento de error. (b)~Un sedimento real no tiene por qué ser normal en unidades phi: una pequeña
mezcla de granos mucho más gruesos o mucho más finos cambia mucho $\varphi_5$ o
$\varphi_{95}$, y poco $\varphi_{16}$ y $\varphi_{84}$, así que la segunda
estimación crece mientras la primera no. Las dos difieren más en las mezclas, como una arena con algo de grava y lodo, o en un depósito bimodal.
\end{solution}

\begin{solution}{exo:b1:clastic-sediments:6}
(a)~$\sin\theta = T/L = 1.5/15 = 0.1$, así que $\theta = \ang{5.74}$. (b)~$a =
c\tan\theta = 0.5 \times 0.1005 = \qty{0.050}{cm}$ por minuto, o
\qty{3.0}{cm} por hora. (c)~Un lecho que se eleva varios centímetros por hora recibe arena que cae de la suspensión muy deprisa: una corriente muy cargada
que se frena, como una crecida en descenso o la cola de una \omterm{def:b1:clastic-sediments:turbidite}{corriente de turbidez}.
\end{solution}

\begin{solution}{exo:b1:clastic-sediments:7}
(a)~$(350 + 20 + 5 + 340 + 15 + 30)/6 = 760/6 = 126.7$: una dirección sudeste, aunque cada lectura queda a menos de \ang{20} del norte. (b)~$S =
\sum\sin\alpha_i = 0.672$, $C = \sum\cos\alpha_i = 5.692$, así que $\bar\alpha =
\operatorname{atan2}(0.672, 5.692) = \ang{6.7}$, y $\bar R =
\sqrt{0.672^2 + 5.692^2}/6 = 5.732/6 = 0.955$. (c)~La media aritmética trata
350 y 005 como muy separados, mientras que sobre el círculo están a \ang{15}: los \omterm{def:b1:spherical-astronomy:horizontal}{acimutes}
deben promediarse como vectores unitarios. El $\bar R$ grande muestra que las lecturas están
agrupadas muy de cerca alrededor de \ang{7}.
\end{solution}

\begin{solution}{exo:b1:clastic-sediments:8}
A tiene más de un \qty{5}{\percent} de arcilla: inmaduro, y también composicionalmente inmaduro, con su \omterm{def:b1:rock-forming-minerals:groups}{feldespato} y sus fragmentos de roca; depositado deprisa
cerca de su fuente, en un abanico aluvial o en una llanura de inundación, o por una \omterm{def:b1:clastic-sediments:turbidite}{corriente de turbidez}. B tiene menos de un \qty{5}{\percent} de arcilla y
está bien clasificado ($\sigma_I < 0.5\,\varphi$), pero sus granos son subangulosos ($\rho < 3$): maduro, por ejemplo una barra de río o una arena de plataforma
cribada por las corrientes. C no tiene arcilla, está muy bien clasificado y sus granos están redondeados ($\rho > 3$): supermaduro, y con un \qty{98}{\percent} de cuarzo
también composicionalmente maduro, como en una playa o en un campo de \omterm{def:b1:clastic-sediments:bedform}{dunas} del desierto.
\end{solution}

\begin{solution}{exo:b1:clastic-sediments:9}
(a)~$1.5/\sqrt{9.81 \times 0.3} = 0.87$; $1.0/\sqrt{9.81 \times 2} = 0.23$.
(b)~$h = U^2/g = 1.5^2/9.81 = \qty{0.23}{m}$; $U = \sqrt{9.81 \times 2} =
\qty{4.43}{m.s^{-1}}$. (c)~El número de Froude crece a medida que disminuye la profundidad a una velocidad dada: las \omterm{def:b1:field-lab-methods-1:thin-section}{láminas delgadas} y rápidas de agua, como el reflujo de una ola
en una playa, alcanzan fácilmente números de Froude cercanos a 1, mientras que un río de varios metros de profundidad necesitaría varios metros por segundo, algo que pocos ríos alcanzan
salvo en crecidas.
\end{solution}

\begin{solution}{exo:b1:clastic-sediments:10}
(a)~Con los puntos medios $-0.5, 0.5, \ldots, 4.5$: $m = (-1 + 4 + 45 + 100 + 52.5
+ 22.5)/100 = 2.23\,\varphi$; las desviaciones cuadráticas ponderadas por los porcentajes suman \num{107.71}, así que $s = \sqrt{1.0771} = 1.04\,\varphi$:
mal clasificada. (b)~El porcentaje acumulado más grueso es 0, 2, 10, 40, 80, 95 y
100 en $-1, 0, \ldots, 5\,\varphi$. Entonces $\varphi_5 = 0 + 3/8 = 0.375$,
$\varphi_{16} = 1 + 6/30 = 1.20$, $\varphi_{50} = 2 + 10/40 = 2.25$,
$\varphi_{84} = 3 + 4/15 = 3.27$, $\varphi_{95} = 4.0$. $M_z = (1.20 + 2.25 +
3.27)/3 = 2.24\,\varphi$ y $\sigma_I = (3.27 - 1.20)/4 + (4.0 - 0.375)/6.6 =
0.517 + 0.549 = 1.07\,\varphi$. (c)~Los dos métodos concuerdan en unos
\num{0.01} en la media y en \num{0.03} en la \omterm{def:b1:clastic-sediments:sorting}{clasificación}: una arena fina
($2.24\,\varphi$ son \qty{0.21}{mm}), mal clasificada.
\end{solution}

\begin{solution}{exo:b1:clastic-sediments:11}
(a)~$U = \sqrt{0.02 \times 9.81 \times 50 \times 0.005/0.003} =
\qty{4.04}{m.s^{-1}}$, o 7.9 nudos. (b)~$(\Delta\rho/\rho)\,h = U^2
C_f/(gS) = 25.7^2 \times 0.003/(9.81 \times 0.005) = \qty{40}{m}$: por ejemplo
una corriente de \qty{400}{m} de espesor con un exceso de densidad del \qty{10}{\percent}, lo que muestra lo grande y densa que debió de ser la corriente de 1929. (c)~La velocidad
crece como $\sqrt S$: $4.04/\sqrt 5 = \qty{1.81}{m.s^{-1}}$. La corriente más lenta tiene menos turbulencia y deposita sus granos más gruesos; se vuelve menos densa, se frena
aún más, y extiende su carga sobre la llanura como un \omterm{def:b1:clastic-sediments:turbidite}{estrato gradado}.
\end{solution}

\begin{solution}{exo:b1:clastic-sediments:12}
(a)~$0.7 \times \qty{5e10}{kg}/\qty{2000}{kg.m^{-3}} = \qty{1.75e7}{m^3}$ al
año. (b)~Repartido sobre un frente de \qty{30}{m} de altura y \qty{100}{km} de anchura, avanza el frente $\num{1.75e7}/(30 \times \num{1e5}) = \qty{5.8}{m}$ al
año. (c)~$a = v\tan\alpha = 5.83 \times \tan\ang{0.5} = \qty{0.051}{m}$, unos
\qty{5}{cm} al año, cien veces menos que el avance; la sucesión de \qty{30}{m} se forma en
$30/0.051 = 590$ años.
\end{solution}

\begin{solution}{pb:b1:clastic-sediments:1}
\textbf{1.} Unidad~1: $(5.0 + 6.6 + 8.4)/3 = 6.67\,\varphi$, \qty{0.0098}{mm}
(\qty{9.8}{\micro m}); unidad~2: $3.47\,\varphi$, \qty{0.090}{mm}; unidad~3:
$1.80\,\varphi$, \qty{0.29}{mm}; unidad~4: $0.73\,\varphi$, \qty{0.60}{mm}.

\textbf{2.} Unidad~1: $(8.4 - 5.0)/4 + (9.6 - 4.2)/6.6 = 0.850 + 0.818 =
1.67\,\varphi$; unidad~2: $0.550 + 0.545 = 1.10$; unidad~3: $0.275 + 0.273 = 0.55$;
unidad~4: $0.725 + 0.712 = 1.44$.

\textbf{3.} Unidad~1: limo fino, mal clasificado; unidad~2: arena muy fina, mal
clasificada; unidad~3: arena media, moderadamente bien clasificada; unidad~4: arena gruesa,
mal clasificada.

\textbf{4.} Unidad~3. En la desembocadura del río la corriente y las olas mueven la arena una y otra vez y se llevan los finos, mientras que el lodo del prodelta se asienta desde la suspensión como una mezcla de todos los tamaños finos que llegan a él.

\textbf{5.} El tamaño $8\,\varphi$ queda a $1.4/1.8$ del camino de $\varphi_{50}$ a $\varphi_{84}$, así que $50 + (1.4/1.8) \times 34 = 76$ por ciento de la
muestra es más gruesa, y unos \qty{24}{\percent} es de tamaño arcilla. Tales partículas finas se asientan solo desde agua casi en reposo: el lugar estaba por debajo del alcance de las olas
y las corrientes.

\textbf{6.} Un canal transporta tanto grava, que rueda por su fondo en las crecidas, como arena; el depósito mezcla el lag de grava del flujo más profundo con la arena de los flujos más bajos, así que cubre muchas clases.

\textbf{7.} La media crece de \qty{0.0098}{mm} a \qty{0.090}{mm} y \qty{0.29}{mm}: la sucesión se engrosa hacia arriba.

\textbf{8.} Unidad~1, lodo laminado: el prodelta; unidad~2, limo y arena alternados: el frente deltaico inferior; unidad~3, arena con \omterm{def:b1:clastic-sediments:cross-bedding}{estratificación cruzada}: las barras de
desembocadura del frente deltaico superior.

\textbf{9.} Cada zona llegó a quedar sobre la que antes estaba delante de ella: el delta avanzó hacia delante, mar adentro, sobre el lugar.

\textbf{10.} Al menos \qty{21}{m}, el espesor de las unidades 1 a 3, ya que la sucesión llenó el agua desde el fondo del prodelta hasta la barra de desembocadura cerca de la superficie;
más, porque los lodos se han compactado desde entonces (\cref{ch:b1:diagenesis-basins}).

\textbf{11.} Un canal distributario de la llanura deltaica, que corta la barra de desembocadura a medida que el delta avanza; sus cantos son su \omterm{def:b1:erosion-transport:load}{carga de fondo}.

\textbf{12.} No: una barra de meandro empieza con una superficie de erosión y una lag, y se vuelve más fina hacia arriba, hasta el lodo de la llanura de inundación. El frente deltaico se engrosa
hacia arriba desde el lodo.

\textbf{13.} $\qty{16}{m}/\num{4000}\text{ años} = \qty{4.0}{mm}$ al año.

\textbf{14.} $v = a/\tan\alpha = 0.0040/\tan\ang{0.5} = 0.0040/0.008727 =
\qty{0.458}{m}$ al año.

\textbf{15.} $0.458 \times 4000 = \qty{1830}{m}$, unos \qty{1.8}{km}.

\textbf{16.} $\qty{21}{m}/\qty{4.0}{mm.yr^{-1}} = 5250$ años.

\textbf{17.} El frente deltaico recibe la mayor parte de la arena cerca de las desembocaduras, el prodelta mucho menos y la llanura deltaica poco, salvo en las crecidas; las
desembocaduras activas también se desplazan de un lugar a otro, así que un punto puede recibir mucho durante un tiempo y luego casi nada. Una sola tasa entre dos estratos datados es un
promedio de estos cambios, y los lodos han perdido parte de su espesor por compactación.

\textbf{18.} La suma es 1544, así que la media aritmética es \ang{128.7}, una dirección sudeste: carece de sentido, pues las lecturas se dispersan alrededor del norte y los valores cercanos a 340 y a 010 se promedian como si estuvieran muy separados.

\textbf{19.} $S = \sum\sin\alpha_i = 1.698$ y $C = \sum\cos\alpha_i =
11.248$.

\textbf{20.} $\bar R = \sqrt{1.698^2 + 11.248^2}/12 = 11.376/12 = 0.948$: las lecturas están
muy agrupadas alrededor de una dirección.

\textbf{21.} Los ríos fluían, y el delta construyó, hacia aproximadamente el norte (\cref{prop:b1:clastic-sediments:azimuths}); la tierra estaba al sur.

\textbf{22.} Porque los foresets se midieron en estratos que ya no están horizontales: sus direcciones de \omterm{def:b1:geological-maps:attitude}{buzamiento} son las de los estratos inclinados más la propia inclinación.
Girar cada foreset alrededor del \omterm{def:b1:geological-maps:attitude}{rumbo} de la estratificación en el ángulo del \omterm{def:b1:geological-maps:attitude}{buzamiento} de esta restaura la estratificación, y con ella los foresets, a su \omterm{def:b1:geological-maps:attitude}{actitud} original.

\textbf{23.} Las corrientes de marea entran y salen, y construyen conjuntos que buzan en dos sentidos opuestos, como en la \omterm{def:b1:clastic-sediments:cross-bedding}{estratificación cruzada} en espina de pescado; el diagrama de rosa tendría
dos lóbulos a unos \ang{180} de distancia. Las doce lecturas forman un solo lóbulo, con $\bar R$ cercano a 1: una corriente de río que fluía en un solo sentido.

\textbf{24.} $\bar\alpha = \operatorname{atan2}(1.698, 11.248) = \ang{8.6}$: el \omterm{def:b1:spherical-astronomy:horizontal}{acimut} medio de la \omterm{def:b1:clastic-sediments:cross-bedding}{paleocorriente} de la unidad~3 es \ang{9}, un poco al este del norte.
\end{solution}
